\documentclass[11pt]{article}
\usepackage{icmlko}
\begin{document}

\title{펀딩은 표를 찾아보고 있었다 --- 점수의 $83\%$ 가 스물한 칸짜리
조회표다}
\author{월드모델 실험실}
\date{노트 471 · 2026년 8월 3일}
\maketitle

\begin{abstract}
노트 470 이 펀딩 점수의 $37\%$ 가 무리 \emph{사이} 순위이고 그 무리가
\texttt{category}(21종)라는 것을 보였다. 그러면 바로 물어야 한다 ---
\textbf{모델 없이 무리만 알면 얼마나 맞히나.}

학습 행에서 무리별 라벨 평균을 내고, 유보 행을 \emph{그 평균으로만}
예보했다. 모델도 축도 없다. 스물한 줄짜리 조회표다.

\begin{center}
\begin{tabular}{lrrrr}
\hline
도메인 & 챔피언 & 무리평균만 & 무리 몫 & 모델이 얹은 것 \\
\hline
\textbf{펀딩} (\texttt{grp} 21) & $+0.3640$ & $\mathbf{+0.3022}$ & $\mathbf{83\%}$ & $\mathbf{+0.0618}$ \\
모바일 (\texttt{gen} 10) & $+0.6223$ & $+0.2854$ & $46\%$ & $+0.3369$ \\
게임 (\texttt{gen} 6) & $+0.6642$ & $+0.2123$ & $32\%$ & $+0.4519$ \\
웹툰 (\texttt{grp} 3) & $+0.4798$ & $+0.1534$ & $32\%$ & $+0.3264$ \\
만화 (\texttt{gen} 13) & $+0.3478$ & $+0.1045$ & $30\%$ & $+0.2433$ \\
애니 (\texttt{gen} 14) & $+0.5731$ & $+0.1311$ & $23\%$ & $+0.4420$ \\
세계애니 (\texttt{gen} 11) & $+0.5616$ & $+0.0972$ & $\mathbf{17\%}$ & $\mathbf{+0.4644}$ \\
\hline
\end{tabular}
\end{center}

\textbf{펀딩에서 모델이 조회표 위에 얹은 것은 $+0.0618$ 뿐이다.}
세계애니는 $+0.4644$ 를 얹는다 --- \textbf{일곱 배 반}.
\end{abstract}

\begin{figure}[t]
\centering
\includegraphics[width=\textwidth]{figs/f1.pdf}
\caption{\textbf{가} --- 회색이 조회표로 되는 부분, 색이 모델이 얹은
부분. \textbf{나} --- 무리 몫. 펀딩만 다른 세상에 있다.}
\end{figure}

\section{왜 펀딩만인가}

기계는 뻔하다. 펀딩의 라벨은 모금액이고, \texttt{category} 는 ``게임''
``출판'' ``공연'' 같은 갈래다. \textbf{갈래마다 전형적인 모금 규모가
자릿수로 다르다.} 그러니 갈래만 알아도 순위의 대부분이 선다.

다른 도메인의 \texttt{gen} 은 그렇지 않다 --- 애니의 ``액션''과 ``드라마''
는 흥행 규모가 자릿수로 다르지 않다. 그래서 무리 몫이 $17 \sim 46\%$ 에
머문다.

\textbf{이것은 펀딩 자료가 나쁘다는 말이 아니다.} 갈래별 규모 차이는
\emph{진짜 정보}다. 다만 \textbf{그 정보는 조회표 한 장에 다 들어간다} ---
2년치 축 서른여섯 개와 나무 서른두 그루가 그 위에 $+0.06$ 을 얹는다.

\section{판을 다시 읽는다}

이 표는 판 $0.4828$ 을 다르게 읽게 한다. 도메인마다 \textbf{모델이 실제로
일한 몫}이 다르다.

\begin{itemize}
\item \textbf{모델이 제일 많이 일하는 곳} --- 세계애니 $+0.4644$ ·
게임 $+0.4519$ · 애니 $+0.4420$.
\item \textbf{제일 적게 일하는 곳} --- \textbf{펀딩 $+0.0618$}. 판
무게가 $529/3{,}430 = 15\%$ 인데, 그 $15\%$ 의 대부분은 조회표가 낸다.
\end{itemize}

\textbf{판에서 높은 도메인이 모델이 많이 일한 곳도 아니다} --- 모바일은
판 $+0.62$ 로 둘째인데 얹은 몫은 $+0.34$ 로 다섯째다.

\section*{조항 --- 도메인 점수 옆에 조회표 값을 적는다}

\begin{quote}
\textbf{조항} --- 도메인 점수를 적을 때 \textbf{무리평균만 예보의 rho} 를
같이 적는다(무리 열쇠가 있는 도메인만). 그 값이 도메인 점수의 절반을
넘으면 \textbf{그 도메인에서 모델이 하는 일은 대부분 조회다} --- 개선을
그 도메인에서 찾지 않는다.
\end{quote}

노트 470 이 ``무리 안 판''을 붙였고 이 노트가 ``무리평균만''을 붙인다.
둘은 짝이다 --- 하나는 \emph{무리를 빼면 얼마 남나}, 다른 하나는
\emph{무리만 있으면 얼마 되나}.

\section*{남은 것}

\begin{itemize}
\item \textbf{펀딩을 판에서 어떻게 다룰까} --- 무게 $15\%$ 를 그대로
두면 판이 조회표를 $15\%$ 만큼 보상한다. 뺄 수는 없다(채점 집합이
바뀌면 2년치와 못 견준다, 노트 378 조건 ④). \textbf{적어 두고 읽을
때 빼는 것}이 지금 할 수 있는 전부다.
\item \textbf{모바일 $46\%$ 도 높다} --- \texttt{gen} 열 무리인데 무엇이
그렇게 갈리나.
\item \textbf{조회표를 이기는 게 목표가 되면} 도메인마다 ``얹은 몫''을
순위로 삼는 판을 따로 만들 수 있다. 지금 판을 대체하는 게 아니라
\emph{옆에} 두는 것.
\end{itemize}

\end{document}
